\begin{ouverture}[calcul]{1}{Les nombres jusqu’au milliard}{Partie 1 · Nombres entiers et calculs}
{Millions · Classes · Comparaisons · Demi-droite}
\objectif{Lire, écrire et décomposer les nombres jusqu’à neuf chiffres.}
\objectif{Comparer, encadrer et intercaler de grands nombres.}
\objectif{Placer et lire des nombres sur une demi-droite graduée.}
\prerequis{Numération jusqu’à 999 999 (classe des mille).}
\end{ouverture}
\pagestyle{otMain}\essentiel\fontsize{11}{15}\selectfont
\phantomsection\label{cours-1-1}
\essSection{1}{La classe des millions}
\begin{definition}Après la classe des unités simples et la classe des mille vient la classe des millions : unités, dizaines et centaines de millions. Le nombre $45\,208\,316$ se lit « quarante-cinq-millions-deux-cent-huit-mille-trois-cent-seize ».\end{definition}\begin{center}\renewcommand{\arraystretch}{1.35}\begin{tabular}{c|ccc}\toprule
Classe & des millions & des mille & des unités\\\midrule
Rangs & c d u & c d u & c d u\\
$45\,208\,316$ & 0 4 5 & 2 0 8 & 3 1 6\\
\bottomrule\end{tabular}\end{center}
\begin{propriete}Chaque rang vaut dix fois le rang placé à sa droite : un million, c’est mille milliers, soit dix centaines de mille. Les zéros gardent la place des rangs vides : $305\,000\,420$ contient 5 unités de millions et aucune dizaine de mille.\end{propriete}
\begin{exemple}$45\,208\,316=45\times1\,000\,000+208\times1\,000+316$. On peut aussi décomposer chiffre par chiffre : $4\times10\,000\,000+5\times1\,000\,000+2\times100\,000+8\times1\,000+3\times100+1\times10+6$.\end{exemple}
\begin{savoirfaire}[Chiffre ou nombre]Dans $45\,208\,316$, le chiffre des dizaines de mille est 0, mais le nombre de milliers est $45\,208$ : tous les paquets de mille contenus. Cette lecture par paquets sert pour les échanges et les ordres de grandeur.\end{savoirfaire}
\clearpage
\phantomsection\label{cours-1-2}
\essSection{2}{Comparer et placer les grands nombres}
\begin{propriete}Pour comparer deux entiers : le moins long est le plus petit ; à longueur égale, comparer rang par rang depuis la gauche. $9\,987\,654<10\,000\,000$ : sept chiffres contre huit.\end{propriete}
\label{fig-cours-1-2-cours-droite-millions}\figureGeo{cours-droite-millions}

\begin{savoirfaire}[Encadrer au rang demandé]Encadrer $7\,802\,540$ entre deux millions consécutifs : $7\,000\,000<7\,802\,540<8\,000\,000$. À la centaine de mille près : $7\,800\,000<7\,802\,540<7\,900\,000$. Le rang demandé fixe les bornes.\end{savoirfaire}
\begin{attention}Les espaces séparent les classes par paquets de trois chiffres, en partant de la droite : ils aident la lecture mais ne changent pas le nombre. Sur une demi-droite, lire la valeur d’un intervalle avant de placer quoi que ce soit.\end{attention}

\clearpage\rubrique{Découverte}
\begin{blocLivre}{ex-1-1}\begin{exercice}[1]{Écrire en chiffres}
Écrire en chiffres : « trois-millions-quarante-mille-deux-cents » ; « soixante-millions-cinq-cent-mille » ; « huit-cent-millions-neuf ».
\end{exercice}\end{blocLivre}
\begin{blocLivre}{ex-1-2}\begin{exercice}[1]{Lire par classes}
Écrire en lettres $45\,208\,316$, puis donner sa classe des mille.
\end{exercice}\end{blocLivre}
\begin{blocLivre}{ex-1-3}\begin{exercice}[1]{La valeur des chiffres}
Dans $305\,000\,420$, donner le chiffre des unités de millions, puis la valeur du chiffre 3.
\end{exercice}\end{blocLivre}
\begin{blocLivre}{ex-1-4}\begin{exercice}[1]{Nombre de milliers}
Pour $45\,208\,316$ : quel est le chiffre des unités de mille ? Combien ce nombre contient-il de milliers entiers ?
\end{exercice}\end{blocLivre}
\begin{blocLivre}{ex-1-5}\begin{exercice}[1]{Comparer}
Compléter par $<$ ou $>$ : $9\,987\,654\ \ldots\ 10\,000\,000$, puis $45\,320\,000\ \ldots\ 45\,230\,000$.
\end{exercice}\end{blocLivre}

\rubrique{Entraînement}
\begin{blocLivre}{ex-1-6}\begin{exercice}[2]{Décomposer et recomposer}
Décomposer $7\,080\,205$ avec les valeurs de chaque chiffre, puis écrire le nombre $6\times10\,000\,000+3\times10\,000+5\times100$.
\end{exercice}\end{blocLivre}
\begin{blocLivre}{ex-1-7}\begin{exercice}[2]{Sur la demi-droite}
Lire les nombres repérés par A, B et C sur la demi-droite graduée.
\label{fig-ex-1-7}\figureGeo{num-7}
\end{exercice}\end{blocLivre}
\begin{blocLivre}{ex-1-8}\begin{exercice}[2]{Encadrer}
Encadrer $7\,802\,540$ entre deux millions consécutifs, puis entre deux dizaines de mille consécutives.
\end{exercice}\end{blocLivre}
\begin{blocLivre}{ex-1-9}\begin{exercice}[2]{Intercaler}
Donner un nombre compris entre $64\,999\,000$ et $65\,001\,000$, dont le chiffre des centaines est 4.
\end{exercice}\end{blocLivre}
\begin{blocLivre}{ex-1-10}\begin{exercice}[2]{Ordonner}
Ranger dans l’ordre décroissant : $9\,900\,000$ ; $10\,100\,000$ ; $9\,090\,000$ ; $10\,010\,000$.
\end{exercice}\end{blocLivre}

\rubrique{Problèmes}
\begin{blocLivre}{ex-1-11}\begin{exercice}[3]{Les populations}
Quatre pays comptent $67\,800\,000$, $47\,400\,000$, $83\,200\,000$ et $10\,500\,000$ habitants. Ranger ces populations dans l’ordre croissant, puis encadrer la plus grande entre deux dizaines de millions consécutives.
\end{exercice}\end{blocLivre}
\begin{blocLivre}{ex-1-12}\begin{exercice}[3]{Des paquets de mille}
Un entrepôt doit expédier $2\,507\,000$ brochures en cartons de $1\,000$. Combien de cartons remplit-il ? Et s’il reçoit encore 500 brochures ?
\end{exercice}\end{blocLivre}
\begin{blocLivre}{ex-1-13}\begin{exercice}[3]{Le compteur du site}
Un site affiche $9\,999\,500$ visites. Combien de visites manquent pour atteindre dix millions ? Quel nombre affichera le compteur 800 visites plus tard ?
\end{exercice}\end{blocLivre}
\begin{blocLivre}{ex-1-14}\begin{exercice}[3]{Avec six cartes}
Avec les cartes 5, 0, 8, 2, 9 et 1 utilisées chacune une fois, écrire le plus grand nombre possible, puis le plus petit nombre à six chiffres.
\end{exercice}\end{blocLivre}
\begin{blocLivre}{ex-1-15}\begin{exercice}[3]{Vrai ou faux ?}
Lina affirme : « $9\,999\,999$ est plus grand que $10\,000\,001$ car il contient plus de 9. » Corriger son raisonnement, puis donner l’écart entre les deux nombres.
\end{exercice}\end{blocLivre}

\clearpage
